\begin{abstract}
Gradient boosted decision trees (GBDTs) are widely used for tabular prediction in
high-stakes domains, and have recently been used in blockchain-facing tasks such
as cryptocurrency fraud detection, phishing-account detection, and smart-contract
risk analysis. In these settings, a model provider may need to convince an
auditor, regulator, or smart-contract-mediated application that a prediction or
model update was produced by a specified GBDT pipeline, without revealing the
private training data, model parameters, or user queries. This makes
zero-knowledge proofs for GBDT training and inference a natural tool for
accountable off-chain machine learning.

Existing approaches for proving GBDT training typically rely on generic ZKP
compilers: they build a certification circuit that checks the forest against the
training data, and then prove the circuit execution. This leads to high prover
cost. A more direct approach is to decompose training into a sequence of
algebraic constraints, which inevitably creates many auxiliary objects. Committing to and
opening these objects can then dominate both proof size and prover time. 
Furthermore, batching
them is also difficult, because these objects are heterogeneous and
have different shapes, sizes, and are committed over different domains.

We present \textsc{Terrae}, a zero-knowledge proof system for quantized GBDT
training and inference based on KZG polynomial commitments. \textsc{Terrae}
avoids both monolithic circuit compilation and isolated auxiliary commitments by
using the structure of GBDT training and batching the resulting proof objects. We
introduce two batching techniques: \emph{domain-lifting batching} for linear
constraints and \emph{interleaving batching} for non-linear constraints. Both
techniques work over heterogeneous objects and reduce many constraints to a
single claim without introducing extra polynomial commitments. We also design a
histogram proof that proves the correctness of converting sample-wise data into
the frequency representation used for split search; this technique may be of
independent interest. Our evaluation shows that, compared with prior approaches,
\textsc{Terrae} significantly reduces proof-generation time while adding only a
small proof-size overhead.
\end{abstract}